\section{Wiles 证明策略概述}
\label{sec:ch10-wiles-strategy}
% ============================================================================

\providecommand{\rhobar}{\bar{\rho}}
\providecommand{\Sel}{\operatorname{Sel}}

Wiles 证明最令人震撼的地方，
不只在于它解决了一个三百多年的猜想，
更在于它把“椭圆曲线模性”这样一句几何断言
彻底翻译成了一个关于 Galois 表示形变与 Hecke 代数的同构问题。
从外部看，
费马大定理似乎与 Hecke 局部化、Selmer 群、完整交环风马牛不相及；
但一旦进入 Wiles 的视角，
这些对象反而比原方程本身更接近答案。

\textbf{我们遇到了什么问题？}
模性定理要求从一条半稳定椭圆曲线 $E/\Q$ 出发，
证明它对应某个新形式 $f$。
直接构造 $f$ 几乎不可能，
因为模形式空间太大，
而椭圆曲线的算术数据又以局部 Galois 条件的形式出现。
Wiles 的突破在于：
既然第\ref{ch:galois}章已经告诉我们“模性”可以写成 Galois 表示同构，
那就不妨直接比较两类表示的 \emph{提升空间}。

\textbf{核心思想是什么？}
先固定一个 residual 表示
\[
  \rhobar\colon G_\Q\to \GL_2(\F_\ell).
\]
一边，考虑所有满足给定局部条件的 $\ell$-进提升，
它们由某个泛形变环 $R$ 参数化；
另一边，考虑所有满足同样局部条件的模形式，
它们在局部 Hecke 代数中编码成一个环 $T$。
如果能证明
\[
  R\cong T,
\]
那么每一个允许的 Galois 提升都自动是模的。
Wiles 的核心定理，
不是“直接找到了对应的新形式”，
而是证明这两个看似来自不同世界的局部环其实同构。

\textbf{引入之后能得到什么？}
一旦 $R=T$，
椭圆曲线的 $\ell$-进表示落在 $R$ 上的一个点就对应于 $T$ 上的一个 Hecke 特征值系统，
也就对应于一个新形式。
这把模性问题从“构造对象”变成了“证明两个变形空间一样大”。
接下来，Langlands--Tunnell 提供起点，
Taylor--Wiles 提供比较环大小的方法，
而 $3$--$5$ 技巧则处理最初 residual 表示可能可约的例外情形。

\begin{figure}[ht]
\centering
\begin{tikzpicture}[>=Stealth, font=\small, node distance=1.2cm]
  \node[draw, rounded corners, fill=orange!14, minimum width=4.8cm, minimum height=0.9cm] (res) {残余表示 $\rhobar_{E,3}$};
  \node[draw, rounded corners, fill=yellow!18, below=of res, minimum width=4.8cm, minimum height=0.9cm] (lt) {Langlands--Tunnell 给出残余模性};
  \node[draw, rounded corners, fill=green!10, below left=1.3cm and 1.8cm of lt, minimum width=3.8cm, minimum height=0.9cm] (r) {形变环 $R$};
  \node[draw, rounded corners, fill=blue!8, below right=1.3cm and 1.8cm of lt, minimum width=3.8cm, minimum height=0.9cm] (t) {Hecke 环 $T$};
  \node[draw, rounded corners, fill=purple!10, below=2.0cm of lt, minimum width=4.2cm, minimum height=0.9cm] (rt) {Taylor--Wiles 证明 $R=T$};
  \node[draw, rounded corners, fill=red!8, below=of rt, minimum width=4.6cm, minimum height=0.9cm] (mod) {所有允许提升都来自模形式};
  \node[draw, rounded corners, fill=gray!14, below=of mod, minimum width=4.8cm, minimum height=0.9cm] (trick) {$3$--$5$ 技巧补足可约例外};

  \draw[->, very thick] (res) -- (lt);
  \draw[->, very thick] (lt) -- (r);
  \draw[->, very thick] (lt) -- (t);
  \draw[->, very thick] (r) -- (rt);
  \draw[->, very thick] (t) -- (rt);
  \draw[->, very thick] (rt) -- (mod);
  \draw[->, very thick] (mod) -- (trick);
\end{tikzpicture}
\caption{Wiles 证明策略的逻辑骨架：先有残余模性，再比较形变环与 Hecke 环，最后用 $R=T$ 把每个允许的提升都判定为模。}
\label{fig:ch10-wiles-flow}
\end{figure}


现在按逻辑顺序概述这条证明链。

\subsection*{第一步：先从一个 residual 表示起步}
对半稳定椭圆曲线 $E/\Q$，取 $\ell=3$，得到
\[
  \rhobar_{E,3}\colon G_\Q\to \GL_2(\F_3).
\]
若它绝对不可约，
则其射影像落在 $\PGL_2(\F_3)\cong S_4$ 的可解部分附近，
足以让 Langlands--Tunnell 定理发挥作用。
这个定理断言：
满足适当奇性与可解像条件的二维 mod $3$ 表示是模的。
于是 Wiles 的 lifting 程序至少有了一个真正的“起跑线”。

\subsection*{第二步：把“所有可能提升”做成一个环}
固定 $\rhobar_{E,3}$ 之后，
我们不只看某一条提升，
而是看所有满足局部条件的提升
\[
  \rho\colon G_\Q\to \GL_2(A),
\]
其中 $A$ 是完备局部环。
这些提升的同构类组成一个形变函子；
Mazur 证明，在温和不可约性条件下它由一个泛环 $R$ 表示。
因此，$R$ 的每一个局部点都对应一条“可能来自椭圆曲线或模形式”的 $3$-进表示。

\subsection*{第三步：模形式一侧出现 Hecke 环 $T$}
把满足相同局部条件的权 $2$ 模形式收集起来，
让 Hecke 算子在对应空间上作用，
再在与 $\rhobar_{E,3}$ 相容的极大理想处局部化并完备化，
得到 Hecke 环 $T$。
第\ref{ch:hecke}章已经告诉我们，
特征形式由 Hecke 特征值系统决定；
因此 $T$ 正是“模的提升”在局部环层面的编码。
由 universal property，存在自然满同态
\[
  R\twoheadrightarrow T.
\]
于是问题变成：这是否其实是同构？

\subsection*{第四步：数值判据与 Taylor--Wiles 打补丁}
Wiles 的数值判据把“是同构吗”转化为“两个局部环的切空间与 congruence ideal 一样大吗”。
形变环的切空间由 Selmer 群控制，
Hecke 环的 congruence ideal 则由模形式之间的同余控制。
Taylor--Wiles 的决定性发明是引入一族辅助素数 $Q_n$，
构造带有额外级别结构的环 $R_{Q_n}$ 与 $T_{Q_n}$，
再通过 inverse limit 与 patching 把问题抬升到一个更自由的极限情形，
在极限上先证明
\[
  R_\infty\cong T_\infty,
\]
然后再下降回原始的最小情形。

\subsection*{补充：Selmer 群如何进入 $R=T$}
CSS 对 Wiles 方法的改编特别强调：
所谓“比较两个局部环大小”，
真正承担数值信息的对象并不是环本身的显式方程，
而是由 Galois 上同调定义出的 Selmer 群。

\begin{definition}[上同调定义的 Selmer 群]
设 $M$ 是一个带 $G_\Q$ 作用的离散 $\ell$-初群，
并在每个素数 $v$ 固定一个局部条件
\[
  L_v\subset H^1(\Q_v,M).
\]
则对应的 Selmer 群定义为
\[
  \Sel_{\mathcal L}(\Q,M)
  =\ker\!\left(H^1(\Q,M)\to \prod_v H^1(\Q_v,M)/L_v\right).
\]
对椭圆曲线 $E/\Q$ 与 $M=E[\ell]$，
通常取 $L_v$ 为 Kummer 映射像，得到经典的 $\ell$-Selmer 群。
\end{definition}

\begin{proposition}[Selmer 群与 Tate--Shafarevich 群]
\label{prop:ch10-selmer-sha}
对椭圆曲线 $E/\Q$ 与素数 $\ell$，存在正合列
\[
  0\to E(\Q)\otimes \Z/\ell\Z
  \to \Sel(\Q,E[\ell])
  \to \text{Tate--Shafarevich 群的 }\ell\text{-挠部分}
  \to 0.
\]
因此 Selmer 群同时测量了有理点秩与 Tate--Shafarevich 群的失败局部--整体原理。
\end{proposition}

\begin{proof}[说明]
Kummer 正合列
\[
  0\to E[\ell]\to E(\bar\Q)\xrightarrow{\ell} E(\bar\Q)\to 0
\]
在全局与局部 Galois 上同调之间给出连接映射。
把所有局部可解条件同时施加到 $H^1(\Q,E[\ell])$ 上，
便得到 Selmer 群；
其商正好记录那些在每个 $\Q_v$ 上都可解、但不来自全局点的扭类，
也就是 Tate--Shafarevich 群的相应部分。
\end{proof}

\begin{theorem}[Wiles 数值判据中的对偶 Selmer 公式]
\label{thm:ch10-wiles-selmer-duality}
设 $f$ 是与极大理想 $\mathfrak m$ 对应的模形式，$T_f$ 是其 $\ell$-进 Galois 晶格，$V_f=T_f\otimes \Q_\ell$。
则在适当的最小局部条件下，Poitou--Tate 对偶给出
\[
  \bigl|\Sel(\Q,V_f/T_f)\bigr|
  =
  \bigl|\Sel(\Q,T_f^*(1)/V_f^*(1))\bigr|,
\]
这里右边是对偶表示的 Selmer 群。
更常见地，人们把它写成两边长度或余维相等。
\end{theorem}

\begin{proof}[说明]
局部 Tate 对偶把每个 $H^1(\Q_v,-)$ 上的允许局部条件与其正交补配对起来；
Poitou--Tate 全局对偶则把这些局部配对拼成一个全局 Euler 特征公式。
在 Wiles 所需的最小情形里，
局部条件被精心选择到恰好自对偶，
于是原始 Selmer 群与对偶 Selmer 群拥有相同的有限大小。
这就是切空间长度可以被精确配平的根本原因。
\end{proof}

\textbf{注。}
形变环 $R$ 的 cotangent space 大小由原始 Selmer 群控制，
而 Hecke 环 $T$ 的 congruence ideal 则由模形式间的同余控制。
Wiles 的数值判据所做的，正是借助上面的对偶性把两边的长度比较变成一个闭合的算术恒等式；
这也解释了为什么 Selmer 群是 $R=T$ 定理里的真正“计量单位”。

\begin{theorem}[Ribet 水平下降]
\label{thm:ribet}
设 $p\nmid N$，
$f\in \Sk_2(\GammaO(Np))$ 是归一化 $p$-new 特征新形式，
$\ell\neq p$，且 residual 表示
\[
  \rhobar_{f,\ell}\colon G_\Q\to \GL_2(\F_\ell)
\]
绝对不可约。
若 $\rhobar_{f,\ell}$ 在 $p$ 处无分歧，
则存在归一化新形式
\[
  g\in \Sk_2(\GammaO(N))
\]
使得
\[
  \rhobar_{g,\ell}\cong \rhobar_{f,\ell}.
\]
\end{theorem}

\begin{proof}[证明思路]
这是第\ref{ch:galois}章的核心定理之一。
其本质是：
级别 $Np$ 与级别 $N$ 的模曲线之间存在两条退化映射，
它们生成 $p$-old 部分；
若 residual 表示在 $p$ 处已经无分歧，
则它在 $J_0(Np)[\mathfrak m]$ 中留下的痕迹无法纯粹待在真正的 $p$-new 部分里，
最终会被 Ihara 引理逼回 old 部分。
而 old 部分本来就来自级别 $N$，
故同一份 residual 表示也来自级别 $N$ 的新形式。
\end{proof}

\subsection*{第五步：$3$--$5$ 技巧处理可约例外}
Wiles 证明里最著名的巧思之一，
就是当 $\rhobar_{E,3}$ 不可用时，
并不是随便换一个素数，
而是精确地换到 $5$。

\begin{proposition}[$3$--$5$ 技巧]
\label{prop:ch10-35-trick}
设 $E/\Q$ 是半稳定椭圆曲线。
若 $\rhobar_{E,3}$ 可约，
则 $\rhobar_{E,5}$ 必不可约。
进一步可构造另一条半稳定椭圆曲线 $A/\Q$，满足
\[
  A[5]\cong E[5],
  \qquad
  \rhobar_{A,3}\ \text{绝对不可约}.
\]
若 $A$ 模，则 $E$ 亦模。
\end{proposition}

\begin{proof}[说明]
若 $E[3]$ 与 $E[5]$ 都可约，
则 $E$ 同时带有有理 $3$-与 $5$-isogeny，
从而给出一条 $15$-isogeny。
Mazur 对有理同源的分类表明，
这在半稳定情形中不可能发生。
因此 $E[5]$ 必不可约。

接着在模曲线 $X(5)$ 上寻找与 $E[5]$ 同构的 $5$-torsion 数据，
可构造一条新曲线 $A$，使得 $A[5]\cong E[5]$ 而 $A[3]$ 不可约。
于是 Langlands--Tunnell 先给出 $A[3]$ 的残余模性，
Wiles 的 lifting 定理再推出 $A$ 模。
最后由 $A[5]\cong E[5]$ 共享同一 residual 表示，
再一次应用 lifting，便得到 $E$ 模。
这正是为什么证明里必须用到 $5$，而不是任意别的素数。
\end{proof}

\textbf{注。}
从今天的视角看，
“先证残余模性，再证所有允许的提升都是模的”
已经是一个标准模板。
但在 Wiles 之前，
真正新颖的正是把这个模板组织成可计算的局部环比较问题。
下一节我们就专门解释这个比较问题里的技术词：形变函子、形变环与 $R=T$。
